% !TEX root = ../thesis.tex
% ==============================================================================
%  APPENDIX: TRAINING HYPERPARAMETERS & ARCHITECTURE (appendices/appendix_hyperparams.tex)
% ==============================================================================
%  This appendix documents the exact hyperparameters, optimizer configurations,
%  and architectural specifications used for training the surrogate EBM and
%  lightweight interpreter networks.
% ==============================================================================

\section{جزئیات ابرپارامترهای آموزش}\label{app:appendix_hyperparams}

در این پیوست، مقادیر دقیق ابرپارامترهای استفاده‌شده برای آموزش شبکه‌های انرژی جانشین (\lr{EBM}) و شبکه‌ی مفسر در چهارچوب \lr{ESCI} ارائه شده است. این تنظیمات حاصل آزمایش‌های تجربی متعدد برای ایجاد تعادل میان پایداری آموزش، جلوگیری از بیش‌برازش و حفظ کارایی محاسباتی بر روی سخت‌افزارهای اقتصادی (نظیر پردازنده‌های گرافیکی \lr{NVIDIA T4}) می‌باشند. 

ساختار معماری و تنظیمات بهینه‌سازی شبکه‌ی جانشین مبتنی بر انرژی در جدول~\ref{tab:app_ebm_hyperparams} خلاصه شده است. در این تنظیمات، استفاده از ۸ لایه‌ی توجه متقاطع و نرخ رهاسازی\enfootnote{dropout} ۱۰ درصد برای جلوگیری از اتکای بیش از حد شبکه به الگوهای نویزی، نقش بسزایی داشته است. همچنین، پس از انجماد مدل انرژی، شبکه‌ی مفسر با پیکربندی شرح‌داده‌شده در جدول~\ref{tab:app_interpreter_hyperparams} آموزش می‌بیند.

\begin{table}[!ht]
    \centering
    \small
    \caption{جزئیات ساختاری و ابرپارامترهای آموزش شبکه‌های جانشین و مفسر}
    \vspace{0.04cm}
    \subtabletopcaptrue % chktex 1
    \renewcommand*{\subcapsize}{\footnotesize}
    \renewcommand*{\subcaplabelfont}{\bfseries}
    \renewcommand*{\subcapfont}{\bfseries}

    \subtable[شبکه‌ی انرژی جانشین\label{tab:app_ebm_hyperparams}]{%
        \resizebox{0.5\textwidth}{!}{%
            \renewcommand{\arraystretch}{1}%
            \setlength{\tabcolsep}{3pt}%
            \setlength{\aboverulesep}{0pt}
            \setlength{\belowrulesep}{0pt}

            \begin{tabular}{lc}
                \toprule
                \rowcolor{gray!15} \textbf{پارامتر} & \textbf{مقدار} \\
                \midrule
                \multicolumn{2}{c}{\raisebox{0.06cm}[0pt][0pt]{\textbf{\textit{پیکربندی معماری}}}} \\
                \arrayrulecolor{gray!20}\midrule[0.1pt]
                مدل پایه‌ی تعبیه‌ساز & \lr{\texttt{all-mpnet-base-v2} (SentenceBERT)} \\
                \arrayrulecolor{gray!20}\midrule[0.1pt]
                نوع پدینگ & قابل یادگیری با مقداردهی تصادفی \\
                \arrayrulecolor{gray!20}\midrule[0.1pt]
                طول بیشینه‌ی جملات ($N_{\text{\lr{max}}}$) & ۲۰ (برای \lr{HC3}) / ۲۴ (برای \lr{WikiBio}) \\
                \arrayrulecolor{gray!20}\midrule[0.1pt]
                بُعد فضای پنهان ($d_{\text{\lr{model}}}$) & $768$ \\
                \arrayrulecolor{gray!20}\midrule[0.1pt]
                تعداد لایه‌های خودتوجهی & $2$ \\
                \arrayrulecolor{gray!20}\midrule[0.1pt]
                تعداد لایه‌های توجه متقاطع & $8$ \\
                \arrayrulecolor{gray!20}\midrule[0.1pt]
                تعداد سرهای مکانیزم توجه & $8$ \\
                \arrayrulecolor{gray!20}\midrule[0.1pt]
                تابع فعال‌سازی & \lr{GELU} \\
                \arrayrulecolor{gray!20}\midrule[0.1pt]
                نرخ رهاسازی & $0.1$ \\
                \arrayrulecolor{gray!20}\midrule[0.1pt]
                تعداد لایه‌های پیش‌خور سر انرژی & $2$ \\
                \arrayrulecolor{gray!20}\midrule[0.1pt]
                ضریب انبساط پنهان در پیش‌خور & $2$ \\
                \arrayrulecolor{black}\midrule
                \multicolumn{2}{c}{\raisebox{0.06cm}[0pt][0pt]{\textbf{\textit{پیکربندی بهینه‌سازی و تابع زیان}}}} \\
                \arrayrulecolor{gray!20}\midrule[0.1pt]
                اندازه‌ی بچ & $16$ \\
                \arrayrulecolor{gray!20}\midrule[0.1pt]
                تعداد ایپاک‌ها & $50$ \\
                \arrayrulecolor{gray!20}\midrule[0.1pt]
                نرخ یادگیری (\lr{LR}) & $3 \times 10^{-5}$ \\
                \arrayrulecolor{gray!20}\midrule[0.1pt]
                زمان‌بند نرخ یادگیری & خطی با ۲۰۰ گام گرم‌کردن \\
                \arrayrulecolor{gray!20}\midrule[0.1pt]
                وزن وابستگی محلی در زیان کل ($\lambda$) & $0.9$ \\
                \arrayrulecolor{gray!20}\midrule[0.1pt]
                دمای \lr{InfoNCE} ($\tau$) & $0.1$ \\
                \arrayrulecolor{gray!20}\midrule[0.1pt]
                حاشیه‌ی امن انرژی & $0.5$ \\
                \arrayrulecolor{black}\bottomrule
            \end{tabular}%
        }%
    }%
    \hfill
    \subtable[شبکه‌ی مفسر\label{tab:app_interpreter_hyperparams}]{%
        \resizebox{0.47\textwidth}{!}{%
            \renewcommand{\arraystretch}{1.05}%
            \setlength{\tabcolsep}{3pt}%
            \setlength{\aboverulesep}{0pt}
            \setlength{\belowrulesep}{0pt}

            \begin{tabular}{lc}
                \toprule
                \rowcolor{gray!15} \vspace{-0.1cm} \raisebox{0.08cm}[0pt][0pt]{\textbf{پارامتر}} & \raisebox{0.08cm}[0pt][0pt]{\textbf{مقدار}} \\
                \midrule
                \multicolumn{2}{>{\columncolor{white}}c}{\raisebox{0.06cm}[0pt][0pt]{\textbf{\textit{پیکربندی معماری و استخراج}}}} \\
                \arrayrulecolor{gray!20}\midrule[0.1pt]
                استراتژی تصویرساز فضای مفهومی & کپی‌شده از \lr{EBM} و آزاد برای آموزش \\
                \arrayrulecolor{gray!20}\midrule[0.1pt]
                بُعد خروجی پیش‌خور نهایی (\lr{MLP}) & $256$ \\
                \arrayrulecolor{gray!20}\midrule[0.1pt]
                روش نرمال‌سازی امتیازات خام & \lr{Sigmoid} (توزیع مستقل برنولی) \\
                \arrayrulecolor{gray!20}\midrule[0.1pt]
                روش اعمال ماسک نرم & ضرب درایه‌به‌درایه \\
                \arrayrulecolor{gray!20}\midrule[0.1pt]
                منظم‌سازی پراکندگی & فعال (جریمه‌ی $0.005$ برای نرم $L_1$) \vspace{0.15cm} \\
                \arrayrulecolor{black}\midrule
                \multicolumn{2}{c}{\raisebox{0.06cm}[0pt][0pt]{\textbf{\textit{پیکربندی بهینه‌سازی}}}} \\
                \arrayrulecolor{gray!20}\midrule[0.1pt]
                تعداد ایپاک‌ها & $15$ \\
                \arrayrulecolor{gray!20}\midrule[0.1pt]
                اندازه‌ی بچ & $16$ \\
                \arrayrulecolor{gray!20}\midrule[0.1pt]
                نرخ یادگیری (\lr{LR}) & $1 \times 10^{-5}$ \\
                \arrayrulecolor{gray!20}\midrule[0.1pt]
                حاشیه‌ی انرژی & $1.0$ \\
                \arrayrulecolor{gray!20}\midrule[0.1pt]
                دمای \lr{InfoNCE} & $0.1$ \vspace{0.1cm} \\
                \arrayrulecolor{black}\midrule
                \multicolumn{2}{c}{\textbf{\textit{پیکربندی آموزش مشترک}}\vspace{0.02cm}} \\
                \arrayrulecolor{gray!20}\midrule[0.1pt]
                تنظیم ظریف \lr{EBM} & فعال \\
                \arrayrulecolor{gray!20}\midrule[0.1pt]
                زمان‌بندی تنظیم ظریف & یک‌بار در هر  $5$ ایپاک \\
                \arrayrulecolor{gray!20}\midrule[0.1pt]
                تعداد گام‌های تنظیم ظریف & $5$ گام در هر بار استعلام \\
                \arrayrulecolor{gray!20}\midrule[0.1pt]
                نرخ یادگیری تنظیم ظریف \lr{EBM} & $5 \times 10^{-6}$ \\
                \arrayrulecolor{gray!20}\midrule[0.1pt]
                حاشیه‌ی انرژی تنظیم ظریف & $0.5$ \vspace{0.15cm} \\
                \arrayrulecolor{black}\bottomrule
            \end{tabular}%
        }%
    }
\end{table}

در مجموع، فرآیند آموزش هر دو فاز به ۳۰ ساعت زمان نیاز داشت. این فرآیند به‌طور کامل بر روی پلتفرم رایگان کگل (\lr{Kaggle}) و با بهره‌گیری از پردازنده‌ی گرافیکی \lr{NVIDIA T4 $\times 2$} اجرا گردید.
